\chapter{Contexto y objetivos}
\label{chap:contexto-objetivos}

\section{Contexto tecnológico}
\label{sec:contexto-tecnologico}

Las aplicaciones backend actuales suelen desplegarse en entornos distribuidos, contenerizados y gestionados por proveedores cloud. Esto permite reducir la administración directa de servidores, pero introduce nuevas responsabilidades: definir redes, permisos, secretos, balanceadores, bases de datos, registros de imágenes, observabilidad y procesos de despliegue.

En este contexto, la \gls{infraestructura-como-codigo} permite describir la infraestructura mediante ficheros versionables, revisables y repetibles. Terraform se utiliza en este proyecto para codificar la red, la base de datos, los servicios de ejecución, los repositorios de imágenes, el balanceador de carga, la configuración de DNS y los permisos necesarios~\cite{TerraformDocs}.

La contenerización proporciona una unidad de despliegue estable entre el entorno local y AWS. La API se empaqueta en una imagen Docker y se ejecuta en ECS Fargate, evitando la gestión directa de servidores. La persistencia se apoya en PostgreSQL gestionado mediante RDS, y el acceso público se concentra en un Application Load Balancer protegido por HTTPS.

El enfoque DevOps aparece en la conexión entre código, validación automática, construcción de imagen, publicación en ECR, despliegue Terraform, migraciones de base de datos y pruebas de humo. El proyecto busca que estas piezas formen un camino verificable desde el repositorio hasta un entorno real.

\section{Objetivo general}
\label{sec:objetivo-general}

Diseñar, implementar, desplegar y operar una plataforma cloud para la gestión de transportes sobre AWS, utilizando infraestructura como código y prácticas DevOps para garantizar reproducibilidad, automatización, seguridad básica y capacidad de observación.

\section{Objetivos específicos}
\label{sec:objetivos-especificos}

\begin{itemize}
  \item Diseñar una arquitectura cloud para una plataforma de gestión de transportes sobre AWS.
  \item Implementar un backend con funcionalidades básicas de transportes, vehículos, conductores, usuarios y eventos.
  \item Definir la infraestructura mediante Terraform para obtener entornos reproducibles.
  \item Contenerizar la aplicación y desplegarla en servicios gestionados de AWS.
  \item Automatizar validación, construcción y despliegue mediante CI/CD.
  \item Incorporar logs, métricas, comprobaciones de salud y trazabilidad operativa.
  \item Aplicar medidas básicas de seguridad en autenticación, autorización, secretos e infraestructura.
  \item Evaluar el sistema mediante pruebas funcionales y análisis de ejecución.
  \item Documentar decisiones, limitaciones y líneas de evolución futura.
\end{itemize}

\section{Criterios de éxito}
\label{sec:criterios-exito}

Los criterios de éxito se han formulado como resultados observables:

\begin{itemize}
  \item La solución compila y las pruebas automatizadas principales se ejecutan correctamente.
  \item La API permite gestionar el flujo operativo básico: autenticación, usuarios, transportes, vehículos, conductores, asignación, cambios de estado y eventos.
  \item La base de datos se gestiona mediante migraciones de Entity Framework Core y no mediante cambios manuales no versionados.
  \item El entorno local puede arrancarse con Docker Compose y configuración externa.
  \item La infraestructura de AWS se define en Terraform y usa estado remoto S3 con bloqueo DynamoDB.
  \item La aplicación se despliega como contenedor en ECS Fargate, detrás de ALB y con persistencia en RDS PostgreSQL.
  \item El entorno \texttt{dev} expone un endpoint HTTPS real en \path{api.dev.transitops.net}.
  \item Las migraciones en AWS se ejecutan como tarea ECS puntual, no como efecto lateral opaco del arranque web.
  \item Los secretos no se almacenan en el repositorio y se consumen desde Secrets Manager.
  \item Existe documentación suficiente para explicar arquitectura, despliegue, validación y limitaciones.
\end{itemize}
